% ============================================================
%  Chapter 1 — Introduction
% ============================================================
\chapter{Introduction}
\label{ch:introduction}

\section{Background}
\label{sec:background}

Groundwater is the quiet backbone of water supply across much of the world.
It irrigates fields when rain fails, fills wells when rivers run low, and in
many semi-arid regions it is simply the only reliable source of fresh water.
That dependence has grown faster than the resource can replenish itself. A
recent global assessment of more than 1,600 aquifer systems found that water
tables are falling in most of them, and in many cases the decline has
accelerated since the start of this century \cite{jasechko2024gwdecline}. The
problem is rarely one of total scarcity. It is one of \emph{knowing where the
water is}---which parts of a landscape store and transmit groundwater, and
which do not---so that extraction can be planned rather than guessed.

Pakistan sits near the sharp end of this problem. The country supports a large
agrarian population on a finite and increasingly stressed water budget, and
the Indus Basin irrigation system that feeds the plains is already operating
under growing demand and a shifting climate \cite{ahmed2021indusbasin}.
Across Punjab, repeated meteorological droughts have begun to translate into
hydrological droughts, with measurable effects on agricultural yield
\cite{waseem2022drought,sarwar2022soandrought}. Urban water security is no
more comfortable: rapid, often unplanned growth in the cities of the north has
pushed municipal supply systems past their design limits
\cite{kamran2024watersecurity}. In this setting groundwater is not a backup.
It is the resource that keeps farms and households running between rains.

Few parts of Punjab feel this more acutely than the \studyarea{}. The plateau
is a dissected upland in the north of the province, largely beyond the reach
of the canal network that irrigates the southern plains. Agriculture here is
overwhelmingly rain-fed, and where surface water is absent the population turns
to dug wells, boreholes, and small rainwater-harvesting structures
\cite{ejaz2023soilwater,hafeez2025rwh}. The climate is semi-arid, rainfall is
unevenly distributed across the year, and large tracts of barren and degraded
land limit natural recharge \cite{khan2024afforestation}. Deciding where to
sink a new well, or where to site a recharge structure, therefore carries real
economic weight for farmers and planners alike---and at present that decision
is made with very little spatial guidance.

\section{Groundwater Potential Mapping}
\label{sec:gwpm}

Groundwater potential mapping (GWPM) is the standard response to this kind of
question. The idea is straightforward: groundwater occurrence is controlled by
a handful of measurable surface and near-surface factors---geology and soil,
slope and drainage, rainfall, vegetation, and so on---so a map that combines
these factors sensibly should highlight the areas most likely to hold usable
groundwater. The factors are derived from remote sensing and digital elevation
data, brought together in a geographic information system (GIS), and merged
into a single potential surface that is usually classified into a few ordinal
zones such as \emph{low}, \emph{medium}, and \emph{high}.

For the better part of two decades the dominant way of combining those factors
has been a knowledge-driven weighted overlay, most often structured through the
analytic hierarchy process (AHP). An expert assigns each factor a weight that
reflects its presumed influence on groundwater, the layers are summed, and the
result is broken into classes. The approach has been applied successfully in a
wide range of settings---the Western Ghats of India \cite{arulbalaji2019gisahp},
the Pravara basin of Maharashtra \cite{das2018pravara}, semi-arid catchments in
India and Saudi Arabia \cite{pande2021ahpmif,mallick2019fuzzyahp}, and
crystalline hard-rock terrain in Morocco \cite{benjmel2020morocco}, among many
others. Its appeal is obvious: it is transparent, it encodes genuine hydrogeological
knowledge, and it works even where no field measurements exist.

That same strength is also its main weakness. The weights are fixed in advance
by judgement, not learned from data, and different experts will reasonably
disagree about them. The method assumes that factors combine in a simple linear,
additive way, which need not hold when, for example, rainfall only matters where
the soil is permeable enough to let it infiltrate. And because the output is a
single deterministic surface, it offers no internal estimate of accuracy and no
straightforward way to ask which factors actually drove the result. Comparative
studies that place weighted-overlay methods alongside statistical and
data-driven alternatives have repeatedly found the data-driven models to be
competitive or better, precisely because they relax these assumptions
\cite{arabameri2018shahroud}.

Machine learning (ML) offers a route around the fixed-weight problem. Rather
than being told how to combine the factors, an ML model is shown examples and
left to learn the combination---including nonlinear interactions and
factor-specific thresholds---directly from the data. Over the past decade ML
has moved from the margins to the mainstream of hydrology and water-resources
research, with applications spanning groundwater-level forecasting, water-quality
assessment, and a broad family of spatial susceptibility problems
\cite{sit2020deeplearning,hai2022gwlreview,zhu2022mlwaterquality}. Tree-based
ensembles in particular---random forests and gradient-boosted trees---have
become a default choice for spatial mapping, both because they handle mixed,
correlated predictors gracefully and because they expose feature-importance
measures that keep the model interpretable. The same machinery that has proved
effective for landslide and flood susceptibility mapping transfers naturally to
groundwater potential, which is structurally the same problem: predicting an
ordinal spatial class from a stack of geo-environmental layers.

\section{Problem Statement}
\label{sec:problem}

The most directly relevant prior study for this work is that of Waqas et al.
\cite{waqas2022chakwal}, who assessed groundwater potential in Chakwal District
using a GIS and remote-sensing weighted overlay built on five thematic layers.
It is a useful and locally grounded study, but it leaves several gaps open.
It covers a single administrative district rather than the plateau as a whole,
so it cannot describe how potential varies across the wider region. It relies
on the manual weighted-overlay method, with all of the subjectivity and the
absence of a quantitative accuracy estimate that this implies. And it uses a
relatively small set of input factors, omitting several variables---soil
texture, a wetness index, surface temperature, and spectral water and
vegetation indices---that later work has shown to be informative.

More broadly, while ML has been applied to closely related problems in the
Pothohar region, including rainfall--runoff simulation across four of its basins
\cite{khan2023rainfallrunoff}, and while recent AHP-based mapping has addressed
rainwater-harvesting potential over the same upland \cite{hafeez2025rwh}, there
is to date no machine-learning groundwater potential map covering the full
five-district plateau at fine resolution. The methodological shift from manual
overlay to learned models, which is by now well established elsewhere, has not
yet been carried through for this study area. That is the gap this thesis sets
out to close.

\section{Research Gap}
\label{sec:gap}

Drawing the preceding observations together, three specific shortcomings define
the gap addressed here:

\begin{itemize}
  \item \textbf{Spatial coverage.} Existing groundwater potential assessments in
        the region are confined to single districts. No consistent,
        full-resolution map spans the entire \studyarea{} (Chakwal, Rawalpindi,
        Attock, Jhelum, and Mianwali).
  \item \textbf{Methodology.} Prior regional mapping relies on knowledge-driven
        weighted overlay or AHP, with subjectively assigned weights, linear
        additive combination, and no quantitative measure of predictive
        performance.
  \item \textbf{Feature set and validation.} Earlier studies use a limited
        number of thematic layers and validate, where they validate at all,
        against secondary data alone. A richer multi-sensor feature set and an
        explicit, evidence-based validation strategy have not been brought
        together for this area.
\end{itemize}

\section{Aim and Objectives}
\label{sec:objectives}

The aim of this thesis is to develop and validate a machine-learning framework
for groundwater potential mapping of the \studyarea{}, extending the
weighted-overlay approach of Waqas et al. \cite{waqas2022chakwal} from a single
district to the full plateau and from a manual method to a learned one. The
specific objectives are:

\begin{enumerate}
  \item To assemble a consistent, nine-feature geo-environmental dataset for the
        five-district plateau at 30~m resolution using Google Earth Engine,
        drawing on Landsat~8, SRTM, CHIRPS, MODIS, and SoilGrids.
  \item To generate groundwater potential reference labels through a transparent,
        knowledge-based weighted overlay, so that the learned models have a
        defensible target grounded in established hydrogeological reasoning.
  \item To train and compare four supervised classifiers---random forest,
        XGBoost, LightGBM, and a soft-voting ensemble---and to select the
        best-performing model on held-out data.
  \item To produce a full-resolution groundwater potential map of the plateau and
        to quantify the spatial distribution of low, medium, and high potential
        zones.
  \item To examine which geo-environmental factors most strongly govern the
        predictions, and to interpret these in hydrogeological terms.
  \item To assess the real-world credibility of the map through an independent,
        spatial-realism validation rather than internal accuracy alone.
\end{enumerate}

\section{Research Questions}
\label{sec:questions}

The objectives translate into four questions that the remainder of the thesis
answers in turn:

\begin{enumerate}
  \item Can a supervised machine-learning model reproduce, and then refine, a
        knowledge-based groundwater potential assessment across the full
        \studyarea{}?
  \item Which of the candidate classifiers performs best, and how large is the
        difference between them in practice?
  \item Which geo-environmental features carry the most weight in the model's
        decisions, and are these consistent with the known hydrogeology of the
        plateau?
  \item Does the resulting map agree with independent evidence---known productive
        sites, the spatial structure of recharge, and findings from comparable
        regional studies---beyond merely reproducing its own training labels?
\end{enumerate}

\section{Scope and Limitations}
\label{sec:scope}

A few boundaries on the work should be stated plainly at the outset. The study
area is limited to the five districts of the \studyarea{}, and all analysis is
carried out at the 30~m native resolution of the SRTM and Landsat inputs. Nine
geo-environmental features are used; subsurface lithology, which would ideally
enter as a tenth layer, was not available in digitized form for the full
plateau and is therefore treated as a known limitation rather than an input.

The most important caveat concerns the training labels. Because field
measurements of groundwater are not available across the whole plateau, the
reference labels are themselves derived from the same satellite features through
a weighted overlay. The machine-learning models are thus trained to reproduce a
knowledge-based product, and their high internal accuracy must be read in that
light: it measures how faithfully the model reproduces the overlay, not how well
it matches ground truth. This is a recognized constraint of data-scarce GWPM,
and rather than gloss over it, the thesis confronts it directly by validating
the final map against \emph{independent} evidence---published well sites,
rain-gauge-based recharge zones, and comparable regional studies---in what is
described here as a spatial-realism validation. The reasoning behind this choice
is developed in Chapters~\ref{ch:methodology} and~\ref{ch:results}.

\section{Significance and Contributions}
\label{sec:significance}

This thesis makes the following contributions:

\begin{itemize}
  \item It delivers the first machine-learning groundwater potential map of the
        full five-district \studyarea{} at 30~m resolution, replacing manual
        weighted overlay with a learned model while retaining hydrogeological
        interpretability.
  \item It expands the feature set from the five layers of the reference study to
        nine, incorporating soil texture, a topographic wetness index, land
        surface temperature, and spectral indices, and quantifies the
        contribution of each.
  \item It provides a systematic comparison of four classifiers on a balanced
        450,000-sample dataset, reporting accuracy, macro $F_1$, Cohen's
        $\kappa$, and ROC-AUC, alongside cross-validated stability.
  \item It introduces an honest, evidence-based validation strategy suited to
        data-scarce settings, demonstrating how a map built on knowledge-based
        labels can nonetheless be checked against independent hydrogeological
        reality.
  \item It produces reusable geospatial outputs---classified and probability
        rasters, and a full set of publication-quality figures---that can support
        district-level groundwater planning by local agencies.
\end{itemize}

\section{Study Area at a Glance}
\label{sec:studyarea_intro}

Figure~\ref{fig:intro_studyarea} locates the \studyarea{} within Pakistan and
shows the five districts that make up the analysis region. A fuller description
of the area's physiography, climate, and hydrogeology is given in
Section~\ref{sec:study_area}.

\begin{figure}[H]
  \centering
  \includegraphics[width=0.95\textwidth]{Fig0_Study_Area}
  \caption{Location of the \studyarea{} study area in northern Punjab, Pakistan,
           showing the five constituent districts---Chakwal, Rawalpindi, Attock,
           Jhelum, and Mianwali---over a shaded-relief base. The plateau covers
           approximately \SI{25000}{\kilo\meter\squared}.}
  \label{fig:intro_studyarea}
\end{figure}

\section{Thesis Organization}
\label{sec:organization}

The remainder of the thesis is organized as follows.
Chapter~\ref{ch:literature} reviews the literature on groundwater potential
mapping, the geo-environmental controls on groundwater, and the move from
knowledge-driven overlay to machine-learning methods, situating the present work
within that progression. Chapter~\ref{ch:methodology} describes the study area in
detail and sets out the full methodology: data collection through Google Earth
Engine, preprocessing, weighted-overlay label generation, construction of the
training dataset, the four classifiers, and the validation design.
Chapter~\ref{ch:results} reports the results---exploratory analysis of the
features, comparative model performance, feature importance, the final
groundwater potential map, and the spatial-realism validation.
Chapter~\ref{ch:discussion} interprets these findings hydrogeologically, compares
them against the reference study and other regional work, and confronts the
label-circularity limitation head-on. Chapter~\ref{ch:conclusion} summarizes the
contributions, notes the limitations, and outlines directions for future
work---chief among them the addition of a lithological layer and the
incorporation of field-measured well data.
